%% kapa: exam=2 round=26 year=2015 period=3 subject=compensation_law subject_ko=감정평가 및 보상법규 seq=1 number=1 points=40
%% source: Q-Net 감정평가사 기출문제 (cst00309) artlSeq=5149601 file=extracted/1008/5149601/2015년 제26회 감정평가사 2차시험 기출문제/3교시(감정평가 및 보상법규).pdf
%% conversion: kordoc 4.24.51 -> tools/gen2.py (auto); review=proofread-gpt-sol; flags=none
\begin{longquestion}{【문제 1】}
「공익사업을 위한 토지 등의 취득 및 보상에 관한 법률」(이하 ‘공익사업법’이라 한다)에 따라 도로확장건설을 위해 사업인정을 받은 A는 해당 지역에 위치한 甲의 토지를 수용하고자 甲과 협의를 시도하였다. A는 甲과 보상액에 관한 협의가 이루어지지 않자 공익사업법상의 절차에 따라 관할 토지수용위원회에 재결을 신청하였다. 그런데 관할 토지수용위원회는 「감정평가에 관한 규칙(국토교통부령)」에 따른 ‘감정평가 실무기준(국토교통부 고시)’과는 다르게 용도지역별 지가변동률이 아닌 이용상황별 지가변동률을 적용한 감정평가사의 감정결과를 채택하여 보상액을 결정하였다. 그 이유로 해당 토지는 이용상황이 지가변동률에 더 큰 영향을 미친다는 것을 들었다. 다음 물음에 답하시오. (40점) \par
1) 甲은 보상액 결정이 ‘감정평가 실무기준(국토교통부 고시)’을 따르지 않았으므로 위법이라고 주장한다. 甲의 주장은 타당한가? (20점) \par
2) 甲은 위 토지수용위원회의 재결에 불복하여 공익사업법에 따라 보상금의 증액을 구하는 소송을 제기하고자 한다. 이 소송의 의의와 그 특수성을 설명하시오. (20점) \par
\end{longquestion}
